\section{Literature Review}\label{sec:literature}

\subsection{Existing LaTeX Templates}

Several \LaTeX\ templates are available for academic writing. Standard classes like \texttt{article} and \texttt{extarticle} provide basic functionality but require extensive configuration for professional output \cite{latexcompanion}. Dedicated journal classes (e.g., \texttt{IEEEtran}, \texttt{elsarticle}) offer journal-specific formatting but may be overly restrictive.

\begin{note}
Unlike journal-specific classes, this template balances customization with professional standards, making it suitable for various publication venues.
\end{note}

\subsection{Required Features for Academic Writing}

Based on a survey of publication requirements, essential template features include \cite{academicwriting2023}:

\begin{itemize}
    \item Two-column layout for efficient space utilization
    \item Automatic cross-referencing of figures, tables, and equations
    \item Citation management with multiple bibliography styles
    \item Code listing support for computational research
    \item Modular document structure for complex projects
\end{itemize}

\subsection{Gap Analysis}

Existing templates often suffer from one or more limitations:

\begin{table}[h]
\centering
\caption{Comparison of Existing Templates}
\label{tab:comparison}
\begin{tabular}{lccc}
\toprule
\textbf{Template} & \textbf{Font Issues} & \textbf{Python Required} & \textbf{Modular} \\
\midrule
rho-class & Yes & Yes & Yes \\
IEEEtran & No & No & Limited \\
elsarticle & No & No & Limited \\
\textbf{Seka (This)} & \textbf{No} & \textbf{No} & \textbf{Yes} \\
\bottomrule
\end{tabular}
\end{table}

This template addresses these gaps by providing a robust, dependency-free solution that maintains professional standards.